\chapter{Overview}
\label{ch:overview}

\bw\ is a one-dimensional, planar, Lagrangian hydrocode for shock and
detonation physics, written in Fortran 2018. It computes plate-impact
problems: a flyer strikes a layered target, and shock and release waves run
through the stack. It began as the reference code for ME~697, Intermediate
Shock and Detonation Theory (Purdue University, Fall 2026); from version 1.4
it is also a home for \emph{open source burn models}: reactive burn models
written down in full, built into a public code, and checked in the open.

This manual says what the code computes and why. It is the source of truth
for the public: where the code and an older document disagree, this manual
follows the code, and says so. Decisions and the studies behind them are
recorded in the design notes (\code{doc/design/}), the input deck in
\code{input/README.md}, and each routine in its module's header; this manual
points to them rather than repeating them.

\section{This version}

The manual is written in three milestones (\code{doc/design/manual.md}).
The first, with \bw\ 1.5.0, completed the two chapters on reactive burn:
\cref{ch:burning-cell}, the burning cell (the closure every burn model
shares), and \cref{ch:burn-models}, the six burn models. With 1.7.0, four of
them are built (HVRB, ThermalCascade, FlashFire and StochasticIgnite) and two
are under development. The chapters on the hydrodynamics and the
materials, and the verification chapter, are outlined and come next; until
then the equation sheets in \code{doc/} describe those parts.

\section{How to read an equation here}

An equation that the code implements carries a \emph{trace tag} under it:
the source file and routine that compute it, and the tests that check it,
for example
\trace{modules/tcas_mod.f95}{tc_rate}{test_tcas}
A script run by \code{make test} checks that every routine and test named in
a tag still exists, so a tag cannot silently point at nothing. Each chapter
names the version of \bw\ its equations were last checked against.
Symbols are listed in \cref{app:notation}; units are \bw's throughout:
mm, \textmu s, g/cm$^3$, km/s, GPa, kJ/g (= (km/s)$^2$) and K.

\tocome{What \bw\ is for, a first run (the class problem, step by step),
the layout of a deck, and what each output file holds (milestone M3).}
